\subsection{切触变换}

在随后的两年里，李似乎放下了变换群理论（尽管他与克莱因保持着极密切的联系，后者于 1872 年发表了他的“纲领”），转而研究切触变换、一阶偏微分方程的积分以及这两个理论之间的关系。我们不必在此详述这些问题的历史，只须提及几点看来在变换群理论的起源中起了重要作用的内容。

切触变换的概念同时推广了点变换与互反极变换。大体上说，一个切触变换\(^{1}\)（在 \(C^{n}\) 中）是开集 \(\Omega\)（流形 \(T^{\prime}(C^{n})\)——即 \(C^{n}\) 的余切向量所成的流形——的开集）到另一个开集 \(\Omega^{\prime}\)（\(T^{\prime}(C^{n})\) 中的）上的一个同构，它把 \(\Omega\) 的典范 1-形式变为 \(\Omega^{\prime}\) 的典范 1-形式。换句话说，若 \((x_{1},\ldots,x_{n},p_{1},\ldots,p_{n})\) 描述 \(T^{\prime}(C^{n})\) 的典范坐标，则切触变换是一个同构 \((x_{i},p_{i})\mapsto(X_{i},P_{i})\)，它满足关系 \(\sum_{i=1}^{n}P_{i}dX_{i}=\sum_{i=1}^{n}p_{i}dx_{i}\)。这样的变换出现在形如

\[
F \left(x _ {1}, x _ {2}, \dots , x _ {n}, \frac {\partial z}{\partial x _ {1}}, \dots , \frac {\partial z}{\partial x _ {n}}\right) = 0\tag{25.2}
\]

李在研究这些问题的过程中，熟悉了对泊松括号

\[
(f, g) = \sum_ {i = 1} ^ {n} \left(\frac {\partial f}{\partial x _ {i}} \frac {\partial g}{\partial p _ {i}} - \frac {\partial g}{\partial x _ {i}} \frac {\partial f}{\partial p _ {i}}\right)\tag{25.3}
\]

以及类型 (25.1) 的微分算子的括号\(^{2}\) \([X,Y]=XY-YX\)的运用；他把泊松括号 (25.3) 解释为一个与 g 相联系的类型 (25.1) 的变换对 f 的作用，并借此机会观察到，泊松括号的雅可比恒等式意味着与 g 和 h 相对应的微分算子的括号与括号 \((g,h)\) 相联系。求使 \((F,g)=0\) 的函数 g——这出现在雅可比积分偏微分方程 (25.2) 的方法中——对李来说变成了求一个使所给方程保持不变的无穷小切触变换。最后，李被引导去研究函数组 \((u_{j})_{1\leq j\leq m}\)（它们是 \(x_{i}\) 与 \(p_{i}\) 的函数），使得诸括号 \((u_{j},u_{k})\) 是诸 \(u_{h}\) 的函数，他称这些集合为“群”（雅可比实质上已经考虑过它们）。

